\section{Deepfake Detection via Adversarial Robustness Training with DCT
Features}\label{deepfake-detection-via-adversarial-robustness-training-with-dct-features}

\subsection{A Comprehensive Technical
Report}\label{a-comprehensive-technical-report}

\textbf{Project:} DDGAN - Deepfake Detection GAN\\
\textbf{Date:} January 2026\\
\textbf{Author:} Rohan Ramesh

\begin{center}\rule{0.5\linewidth}{0.5pt}\end{center}

\subsection{Table of Contents}\label{table-of-contents}

\begin{enumerate}
\def\labelenumi{\arabic{enumi}.}
\tightlist
\item
  \hyperref[1-executive-summary]{Executive Summary}
\item
  \hyperref[2-introduction]{Introduction}
\item
  \hyperref[3-methodology]{Methodology}
\item
  \hyperref[4-system-architecture]{System Architecture}

  \begin{itemize}
  \tightlist
  \item
    4.1 \hyperref[41-dct-feature-extractor]{DCT Feature Extractor}
  \item
    4.2 \hyperref[42-discriminator-architecture]{Discriminator
    Architecture}
  \item
    4.3 \hyperref[43-generator-architecture]{Generator Architecture}
  \end{itemize}
\item
  \hyperref[5-training-framework]{Training Framework}

  \begin{itemize}
  \tightlist
  \item
    5.1 \hyperref[51-training-philosophy]{Training Philosophy}
  \item
    5.2 \hyperref[52-data-pipeline]{Data Pipeline}
  \item
    5.3 \hyperref[53-training-algorithm]{Training Algorithm}
  \item
    5.4 \hyperref[54-loss-functions]{Loss Functions}
  \item
    5.5 \hyperref[55-optimization-strategy]{Optimization Strategy}
  \end{itemize}
\item
  \hyperref[6-implementation-details]{Implementation Details}
\item
  \hyperref[7-expected-results]{Expected Results}
\item
  \hyperref[8-conclusion]{Conclusion}
\item
  \hyperref[9-references]{References}
\end{enumerate}

\begin{center}\rule{0.5\linewidth}{0.5pt}\end{center}

\subsection{1. Executive Summary}\label{executive-summary}

This report presents a novel deepfake detection system that leverages
\textbf{adversarial robustness training} with \textbf{Discrete Cosine
Transform (DCT)} frequency-domain features. Unlike traditional deepfake
detectors that rely solely on spatial features, our approach:

\begin{itemize}
\tightlist
\item
  \textbf{Extracts frequency-domain features} using 2D DCT, capturing
  manipulation artifacts invisible in the spatial domain
\item
  \textbf{Employs adversarial training} where a generator creates
  perturbations to stress-test the discriminator
\item
  \textbf{Achieves robustness} by training the discriminator to
  correctly classify images even under adversarial attack
\end{itemize}

\textbf{Key Statistics:} - Total Parameters: \textbf{41.8 million} -
Discriminator: \textasciitilde27.8M (ConvNeXt-Tiny backbone) -
Generator: \textasciitilde14M (U-Net architecture) - Input Resolution:
224×224 RGB - Dataset: Celeb-DF v2 (\textasciitilde537K training,
\textasciitilde297K validation images)

\begin{center}\rule{0.5\linewidth}{0.5pt}\end{center}

\subsection{2. Introduction}\label{introduction}

\subsubsection{2.1 Problem Statement}\label{problem-statement}

Deepfake technology has advanced rapidly, enabling the creation of
highly realistic synthetic media. Detecting such manipulations is
critical for: - Preventing misinformation and fraud - Protecting
individual privacy and reputation - Maintaining trust in digital media

\subsubsection{2.2 Motivation for DCT-based
Detection}\label{motivation-for-dct-based-detection}

Traditional CNN-based detectors often fail against: - Post-processing
(compression, resizing) - Adversarial perturbations - Novel generation
techniques

\textbf{Why DCT Features?} 1. \textbf{Frequency artifacts are more
persistent} - GAN-generated images exhibit distinct patterns in the
frequency domain that survive common post-processing 2. \textbf{JPEG
compression signatures} differ between real and synthetic images 3.
\textbf{High-frequency anomalies} that GANs struggle to replicate are
clearly visible in DCT space 4. \textbf{Compact representation} of
manipulation signatures

\subsubsection{2.3 Novel Contributions}\label{novel-contributions}

\begin{enumerate}
\def\labelenumi{\arabic{enumi}.}
\tightlist
\item
  \textbf{DCT-based feature extraction} integrated directly into the
  neural network (trainable end-to-end)
\item
  \textbf{Adversarial robustness training} instead of traditional GAN
  objectives
\item
  \textbf{Frequency-aware generator} with implicit frequency filtering
  via large-kernel convolutions
\item
  \textbf{Transfer learning} from ImageNet adapted for single-channel
  DCT features
\end{enumerate}

\begin{center}\rule{0.5\linewidth}{0.5pt}\end{center}

\subsection{3. Methodology}\label{methodology}

\subsubsection{3.1 Overview}\label{overview}

Our system consists of two primary components trained adversarially:

\begin{verbatim}
┌─────────────────────────────────────────────────────────────────┐
│                    TRAINING FRAMEWORK                           │
├─────────────────────────────────────────────────────────────────┤
│                                                                 │
│   ┌─────────────┐    perturbation    ┌─────────────────────┐   │
│   │  Generator  │ ─────────────────► │   Adversarial       │   │
│   │   (U-Net)   │                    │      Image          │   │
│   └─────────────┘                    └──────────┬──────────┘   │
│         ▲                                       │              │
│         │                                       │              │
│         │ feedback                              ▼              │
│         │                            ┌─────────────────────┐   │
│   Real Image ──────────────────────► │   Discriminator     │   │
│                                      │  (DCT + ConvNeXt)   │   │
│   Fake Image ──────────────────────► │                     │   │
│                                      └─────────────────────┘   │
│                                               │                │
│                                               ▼                │
│                                        Real/Fake Output        │
└─────────────────────────────────────────────────────────────────┘
\end{verbatim}

\subsubsection{3.2 Key Insight: Adversarial Robustness vs Traditional
GAN}\label{key-insight-adversarial-robustness-vs-traditional-gan}

\begin{longtable}[]{@{}
  >{\raggedright\arraybackslash}p{(\linewidth - 4\tabcolsep) * \real{0.2105}}
  >{\raggedright\arraybackslash}p{(\linewidth - 4\tabcolsep) * \real{0.4211}}
  >{\raggedright\arraybackslash}p{(\linewidth - 4\tabcolsep) * \real{0.3684}}@{}}
\toprule\noalign{}
\begin{minipage}[b]{\linewidth}\raggedright
Aspect
\end{minipage} & \begin{minipage}[b]{\linewidth}\raggedright
Traditional GAN
\end{minipage} & \begin{minipage}[b]{\linewidth}\raggedright
Our Approach
\end{minipage} \\
\midrule\noalign{}
\endhead
\bottomrule\noalign{}
\endlastfoot
Generator Goal & Create realistic fakes & Create adversarial
perturbations \\
Discriminator Goal & Detect fakes & Be robust to perturbations \\
Training Target & Nash equilibrium (D confused) & Strong D despite
attacks \\
Final Product & Generator for synthesis & Robust discriminator for
detection \\
\end{longtable}

\begin{center}\rule{0.5\linewidth}{0.5pt}\end{center}

\subsection{4. System Architecture}\label{system-architecture}

\subsubsection{4.1 DCT Feature Extractor}\label{dct-feature-extractor}

The DCT Feature Extractor transforms RGB images into frequency-domain
representations.

\paragraph{4.1.1 Mathematical Foundation}\label{mathematical-foundation}

\textbf{DCT-II Transform:}

The 2D Discrete Cosine Transform is defined as:

\[F(u,v) = \alpha(u)\alpha(v) \sum_{x=0}^{N-1} \sum_{y=0}^{N-1} f(x,y) \cos\left[\frac{\pi(2x+1)u}{2N}\right] \cos\left[\frac{\pi(2y+1)v}{2N}\right]\]

Where: - \(\alpha(0) = \sqrt{1/N}\) (DC component normalization) -
\(\alpha(k) = \sqrt{2/N}\) for \(k > 0\) (AC components)

\textbf{Separable Implementation:}

For computational efficiency, we implement 2D DCT as two 1D transforms:

\[F = D \cdot f \cdot D^T\]

Where \(D\) is the precomputed DCT transformation matrix.

\paragraph{4.1.2 Processing Pipeline}\label{processing-pipeline}

\begin{verbatim}
┌─────────────────────────────────────────────────────────────────┐
│               DCT FEATURE EXTRACTION PIPELINE                    │
├─────────────────────────────────────────────────────────────────┤
│                                                                 │
│   RGB Image [B, 3, 224, 224]                                    │
│        │                                                        │
│        ▼                                                        │
│   ┌─────────────────────────────────────────────────────────┐   │
│   │  Grayscale Conversion (ITU-R BT.601)                    │   │
│   │  Y = 0.299R + 0.587G + 0.114B                           │   │
│   └─────────────────────────────────────────────────────────┘   │
│        │                                                        │
│        ▼                                                        │
│   Grayscale [B, 1, 224, 224]                                    │
│        │                                                        │
│        ▼                                                        │
│   ┌─────────────────────────────────────────────────────────┐   │
│   │  Separable 2D DCT Transform                             │   │
│   │  Step 1: DCT along rows (right multiply with D^T)       │   │
│   │  Step 2: DCT along columns (left multiply with D)       │   │
│   └─────────────────────────────────────────────────────────┘   │
│        │                                                        │
│        ▼                                                        │
│   DCT Coefficients [B, 1, 224, 224]                             │
│        │                                                        │
│        ▼                                                        │
│   ┌─────────────────────────────────────────────────────────┐   │
│   │  Log Scaling (Dynamic Range Compression)                │   │
│   │  output = log(|DCT| + ε), where ε = 1e-6               │   │
│   └─────────────────────────────────────────────────────────┘   │
│        │                                                        │
│        ▼                                                        │
│   DCT Features [B, 1, 224, 224]                                 │
│                                                                 │
└─────────────────────────────────────────────────────────────────┘
\end{verbatim}

\paragraph{4.1.3 Implementation Details}\label{implementation-details}

\begin{Shaded}
\begin{Highlighting}[]
\KeywordTok{class}\NormalTok{ DCT2D(nn.Module):}
    \KeywordTok{def} \FunctionTok{\_\_init\_\_}\NormalTok{(}\VariableTok{self}\NormalTok{, size}\OperatorTok{=}\DecValTok{224}\NormalTok{):}
        \BuiltInTok{super}\NormalTok{().}\FunctionTok{\_\_init\_\_}\NormalTok{()}
        \CommentTok{\# Precompute DCT basis matrix (224×224)}
\NormalTok{        dct\_matrix }\OperatorTok{=} \VariableTok{self}\NormalTok{.\_create\_dct\_matrix(size)}
        \CommentTok{\# Register as buffer (non{-}trainable, auto GPU transfer)}
        \VariableTok{self}\NormalTok{.register\_buffer(}\StringTok{\textquotesingle{}dct\_matrix\textquotesingle{}}\NormalTok{, dct\_matrix)}
    
    \KeywordTok{def}\NormalTok{ forward(}\VariableTok{self}\NormalTok{, x):}
        \CommentTok{\# RGB → Grayscale}
\NormalTok{        gray }\OperatorTok{=} \VariableTok{self}\NormalTok{.rgb\_to\_grayscale(x)}
        
        \CommentTok{\# Separable 2D DCT}
\NormalTok{        dct\_rows }\OperatorTok{=}\NormalTok{ torch.matmul(gray, }\VariableTok{self}\NormalTok{.dct\_matrix.t())}
\NormalTok{        dct\_2d }\OperatorTok{=}\NormalTok{ torch.matmul(}\VariableTok{self}\NormalTok{.dct\_matrix, dct\_rows)}
        
        \CommentTok{\# Clamp for numerical stability (fp16)}
\NormalTok{        dct\_2d }\OperatorTok{=}\NormalTok{ torch.clamp(dct\_2d, }\OperatorTok{{-}}\FloatTok{1e6}\NormalTok{, }\FloatTok{1e6}\NormalTok{)}
        
        \CommentTok{\# Log scaling}
        \ControlFlowTok{return}\NormalTok{ torch.log(torch.}\BuiltInTok{abs}\NormalTok{(dct\_2d) }\OperatorTok{+} \FloatTok{1e{-}6}\NormalTok{)}
\end{Highlighting}
\end{Shaded}

\textbf{Why Log Scaling?} - DCT coefficients span huge dynamic range (DC
component \textgreater\textgreater{} AC components) - Logarithmic
compression normalizes the feature distribution - Prevents gradient
domination by DC component - Aligns with human perceptual logarithmic
response

\begin{center}\rule{0.5\linewidth}{0.5pt}\end{center}

\subsubsection{4.2 Discriminator
Architecture}\label{discriminator-architecture}

The Discriminator classifies images as real or fake using DCT features
processed through a ConvNeXt-Tiny backbone.

\paragraph{4.2.1 Architecture Overview}\label{architecture-overview}

\begin{verbatim}
┌─────────────────────────────────────────────────────────────────────────┐
│                      DISCRIMINATOR ARCHITECTURE                          │
│                        Parameters: ~27.8 Million                         │
├─────────────────────────────────────────────────────────────────────────┤
│                                                                         │
│   INPUT: RGB Image [B, 3, 224, 224]                                     │
│        │                                                                │
│        ▼                                                                │
│   ┌─────────────────────────────────────────────────────────────────┐   │
│   │                    DCT Feature Extractor                        │   │
│   │                    (Non-trainable buffers)                      │   │
│   └─────────────────────────────────────────────────────────────────┘   │
│        │                                                                │
│        ▼                                                                │
│   DCT Features [B, 1, 224, 224]                                         │
│        │                                                                │
│        ▼                                                                │
│   ┌─────────────────────────────────────────────────────────────────┐   │
│   │                    ConvNeXt-Tiny Backbone                       │   │
│   │                 (ImageNet Pretrained, Modified)                 │   │
│   ├─────────────────────────────────────────────────────────────────┤   │
│   │                                                                 │   │
│   │   STEM LAYER (Modified for 1-channel input)                     │   │
│   │   Conv2d(1→96, k=4, s=4) + LayerNorm                           │   │
│   │   Output: [B, 96, 56, 56]                                       │   │
│   │                                                                 │   │
│   │   ─────────────────────────────────────────────                 │   │
│   │                                                                 │   │
│   │   STAGE 1: 3× ConvNeXt Blocks                                   │   │
│   │   Output: [B, 96, 56, 56]                                       │   │
│   │                                                                 │   │
│   │   Downsample: LayerNorm + Conv2d(96→192, k=2, s=2)             │   │
│   │                                                                 │   │
│   │   ─────────────────────────────────────────────                 │   │
│   │                                                                 │   │
│   │   STAGE 2: 3× ConvNeXt Blocks                                   │   │
│   │   Output: [B, 192, 28, 28]                                      │   │
│   │                                                                 │   │
│   │   Downsample: LayerNorm + Conv2d(192→384, k=2, s=2)            │   │
│   │                                                                 │   │
│   │   ─────────────────────────────────────────────                 │   │
│   │                                                                 │   │
│   │   STAGE 3: 9× ConvNeXt Blocks                                   │   │
│   │   Output: [B, 384, 14, 14]                                      │   │
│   │                                                                 │   │
│   │   Downsample: LayerNorm + Conv2d(384→768, k=2, s=2)            │   │
│   │                                                                 │   │
│   │   ─────────────────────────────────────────────                 │   │
│   │                                                                 │   │
│   │   STAGE 4: 3× ConvNeXt Blocks                                   │   │
│   │   Output: [B, 768, 7, 7]                                        │   │
│   │                                                                 │   │
│   └─────────────────────────────────────────────────────────────────┘   │
│        │                                                                │
│        ▼                                                                │
│   ┌─────────────────────────────────────────────────────────────────┐   │
│   │              Global Average Pooling                             │   │
│   │              AdaptiveAvgPool2d(1)                               │   │
│   └─────────────────────────────────────────────────────────────────┘   │
│        │                                                                │
│        ▼                                                                │
│   Feature Vector [B, 768]                                               │
│        │                                                                │
│        ▼                                                                │
│   ┌─────────────────────────────────────────────────────────────────┐   │
│   │              Classification Head                                │   │
│   │              LayerNorm(768) → Linear(768→1)                     │   │
│   └─────────────────────────────────────────────────────────────────┘   │
│        │                                                                │
│        ▼                                                                │
│   OUTPUT: Logit [B, 1] (Real/Fake probability before sigmoid)           │
│                                                                         │
└─────────────────────────────────────────────────────────────────────────┘
\end{verbatim}

\paragraph{4.2.2 ConvNeXt Block
Structure}\label{convnext-block-structure}

Each ConvNeXt block follows this architecture:

\begin{verbatim}
Input [B, C, H, W]
      │
      ├────────────────────────────────┐
      │                                │ (Residual)
      ▼                                │
┌─────────────────────────────────┐    │
│ Depthwise Conv 7×7              │    │
│ (groups=C, padding=3)           │    │
└─────────────────────────────────┘    │
      │                                │
      ▼                                │
┌─────────────────────────────────┐    │
│ LayerNorm                       │    │
└─────────────────────────────────┘    │
      │                                │
      ▼                                │
┌─────────────────────────────────┐    │
│ Pointwise Conv 1×1 (C→4C)       │    │
│ GELU Activation                 │    │
└─────────────────────────────────┘    │
      │                                │
      ▼                                │
┌─────────────────────────────────┐    │
│ Pointwise Conv 1×1 (4C→C)       │    │
└─────────────────────────────────┘    │
      │                                │
      ▼                                │
      + ◄──────────────────────────────┘
      │
      ▼
Output [B, C, H, W]
\end{verbatim}

\paragraph{4.2.3 Input Adaptation
Strategy}\label{input-adaptation-strategy}

The original ConvNeXt expects 3-channel RGB input. We adapt it for
1-channel DCT:

\begin{Shaded}
\begin{Highlighting}[]
\CommentTok{\# Average pretrained RGB weights for single{-}channel initialization}
\NormalTok{original\_conv }\OperatorTok{=}\NormalTok{ backbone.stem[}\DecValTok{0}\NormalTok{]  }\CommentTok{\# Conv2d(3, 96, k=4, s=4)}
\NormalTok{new\_conv }\OperatorTok{=}\NormalTok{ Conv2d(}\DecValTok{1}\NormalTok{, }\DecValTok{96}\NormalTok{, kernel\_size}\OperatorTok{=}\DecValTok{4}\NormalTok{, stride}\OperatorTok{=}\DecValTok{4}\NormalTok{)}

\CommentTok{\# Initialize: average RGB weights}
\NormalTok{new\_conv.weight }\OperatorTok{=}\NormalTok{ original\_conv.weight.mean(dim}\OperatorTok{=}\DecValTok{1}\NormalTok{, keepdim}\OperatorTok{=}\VariableTok{True}\NormalTok{)}
\end{Highlighting}
\end{Shaded}

\textbf{Why this approach?} - Preserves learned low-level feature
detectors from ImageNet - Grayscale features share structure with RGB
luminance channel - Faster convergence vs random initialization

\begin{center}\rule{0.5\linewidth}{0.5pt}\end{center}

\subsubsection{4.3 Generator Architecture}\label{generator-architecture}

The Generator creates adversarial perturbations using a U-Net
architecture with a frequency-aware bottleneck.

\paragraph{4.3.1 Architecture Overview}\label{architecture-overview-1}

\begin{verbatim}
┌─────────────────────────────────────────────────────────────────────────────┐
│                        GENERATOR ARCHITECTURE (U-Net)                        │
│                           Parameters: ~14 Million                            │
├─────────────────────────────────────────────────────────────────────────────┤
│                                                                             │
│   INPUT: RGB Image [B, 3, 224, 224]                                         │
│        │                                                                    │
│   ═════╪═══════════════════════════════════════════════════════════════     │
│   ║    │               ENCODER (Downsampling Path)                    ║     │
│   ║    ▼                                                              ║     │
│   ║   ┌─────────────────┐                                             ║     │
│   ║   │   Enc Block 1   │ Conv3×3→BN→GELU→Conv3×3→BN→GELU             ║     │
│   ║   │   [3 → 64]      │───────────────────────────────┐ e1          ║     │
│   ║   └────────┬────────┘                               │             ║     │
│   ║            │ MaxPool2d(2)                           │             ║     │
│   ║            ▼                                        │             ║     │
│   ║   ┌─────────────────┐                               │             ║     │
│   ║   │   Enc Block 2   │ [B, 64, 112, 112]             │             ║     │
│   ║   │   [64 → 128]    │───────────────────┐ e2        │             ║     │
│   ║   └────────┬────────┘                   │           │             ║     │
│   ║            │ MaxPool2d(2)               │           │             ║     │
│   ║            ▼                            │           │             ║     │
│   ║   ┌─────────────────┐                   │           │             ║     │
│   ║   │   Enc Block 3   │ [B, 128, 56, 56]  │           │             ║     │
│   ║   │   [128 → 256]   │───────────┐ e3    │           │             ║     │
│   ║   └────────┬────────┘           │       │           │             ║     │
│   ║            │ MaxPool2d(2)       │       │           │             ║     │
│   ║            ▼                    │       │           │             ║     │
│   ║   ┌─────────────────┐           │       │           │             ║     │
│   ║   │   Enc Block 4   │ [B, 256, 28, 28]  │           │             ║     │
│   ║   │   [256 → 512]   │───┐ e4    │       │           │             ║     │
│   ║   └────────┬────────┘   │       │       │           │             ║     │
│   ║            │ MaxPool2d(2)       │       │           │             ║     │
│   ║            ▼            │       │       │           │             ║     │
│   ═════════════╪════════════╪═══════╪═══════╪═══════════╪═════════════      │
│                │            │       │       │           │                   │
│   ┌────────────┴────────────────────────────────────────────────────┐       │
│   │                    BOTTLENECK (512 channels)                    │       │
│   │                     [B, 512, 14, 14]                            │       │
│   ├─────────────────────────────────────────────────────────────────┤       │
│   │  ┌───────────────────────────────────────────────────────────┐  │       │
│   │  │ Conv2d(512→512, 3×3) → BatchNorm → GELU                   │  │       │
│   │  └───────────────────────────────────────────────────────────┘  │       │
│   │                          │                                      │       │
│   │                          ▼                                      │       │
│   │  ┌───────────────────────────────────────────────────────────┐  │       │
│   │  │          FrequencyAwareBottleneck (See 4.3.2)             │  │       │
│   │  └───────────────────────────────────────────────────────────┘  │       │
│   │                          │                                      │       │
│   │                          ▼                                      │       │
│   │  ┌───────────────────────────────────────────────────────────┐  │       │
│   │  │ Conv2d(512→512, 3×3) → BatchNorm → GELU                   │  │       │
│   │  └───────────────────────────────────────────────────────────┘  │       │
│   └─────────────────────────────────────────────────────────────────┘       │
│                │            │       │       │           │                   │
│   ═════════════╪════════════╪═══════╪═══════╪═══════════╪═════════════      │
│   ║            │            │       │       │           │             ║     │
│   ║    DECODER │(Upsampling │Path + │Skip   │Connections)             ║     │
│   ║            ▼            │       │       │           │             ║     │
│   ║   ┌─────────────────┐   │       │       │           │             ║     │
│   ║   │ ConvTranspose2d │   │       │       │           │             ║     │
│   ║   │   (Upsample)    │◄──┘       │       │           │             ║     │
│   ║   │ [B, 512, 28,28] │           │       │           │             ║     │
│   ║   └────────┬────────┘           │       │           │             ║     │
│   ║            │ Concat(e4) ────────┘       │           │             ║     │
│   ║            ▼ [B, 1024, 28, 28]          │           │             ║     │
│   ║   ┌─────────────────┐                   │           │             ║     │
│   ║   │   Dec Block 4   │                   │           │             ║     │
│   ║   │ [1024 → 256]    │                   │           │             ║     │
│   ║   └────────┬────────┘                   │           │             ║     │
│   ║            │ Upsample + Concat(e3) ─────┘           │             ║     │
│   ║            ▼ [B, 512, 56, 56]                       │             ║     │
│   ║   ┌─────────────────┐                               │             ║     │
│   ║   │   Dec Block 3   │                               │             ║     │
│   ║   │ [512 → 128]     │                               │             ║     │
│   ║   └────────┬────────┘                               │             ║     │
│   ║            │ Upsample + Concat(e2) ─────────────────┘             ║     │
│   ║            ▼ [B, 256, 112, 112]                                   ║     │
│   ║   ┌─────────────────┐                                             ║     │
│   ║   │   Dec Block 2   │                                             ║     │
│   ║   │ [256 → 64]      │                                             ║     │
│   ║   └────────┬────────┘                                             ║     │
│   ║            │ Upsample + Concat(e1) ───────────────────────────────┘     │
│   ║            ▼ [B, 128, 224, 224]                                   ║     │
│   ║   ┌─────────────────┐                                             ║     │
│   ║   │   Dec Block 1   │                                             ║     │
│   ║   │ [128 → 64]      │                                             ║     │
│   ║   └────────┬────────┘                                             ║     │
│   ═════════════╪══════════════════════════════════════════════════════      │
│                │                                                            │
│                ▼                                                            │
│   ┌─────────────────────────────────────────────────────────────────┐       │
│   │              OUTPUT LAYER                                       │       │
│   │  Conv2d(64→3, 1×1) → Tanh                                       │       │
│   │  Perturbation: [B, 3, 224, 224] ∈ [-1, 1]                       │       │
│   └─────────────────────────────────────────────────────────────────┘       │
│                │                                                            │
│                ▼                                                            │
│   ┌─────────────────────────────────────────────────────────────────┐       │
│   │         ADVERSARIAL IMAGE GENERATION                            │       │
│   │  scaled_perturbation = ε × perturbation (ε = 0.03)              │       │
│   │  adversarial = input + scaled_perturbation                      │       │
│   │  adversarial = clamp(adversarial, valid_range)                  │       │
│   └─────────────────────────────────────────────────────────────────┘       │
│                │                                                            │
│                ▼                                                            │
│   OUTPUT: (adversarial_image, perturbation)                                 │
│                                                                             │
└─────────────────────────────────────────────────────────────────────────────┘
\end{verbatim}

\paragraph{4.3.2 Frequency-Aware
Bottleneck}\label{frequency-aware-bottleneck}

The bottleneck module is specifically designed to learn frequency-domain
perturbations:

\begin{verbatim}
┌─────────────────────────────────────────────────────────────────────────┐
│                    FREQUENCY-AWARE BOTTLENECK                            │
├─────────────────────────────────────────────────────────────────────────┤
│                                                                         │
│   Input [B, 512, H, W]                                                  │
│        │                                                                │
│        ├────────────────────────────────────────────┐ (Residual)        │
│        │                                            │                   │
│        ▼                                            │                   │
│   ┌─────────────────────────────────────────────┐   │                   │
│   │  Depthwise Conv 7×7 (groups=512)            │   │                   │
│   │  - Large kernel captures frequency patterns │   │                   │
│   │  - Efficient: O(512×7×7) vs O(512²×3×3)     │   │                   │
│   └─────────────────────────────────────────────┘   │                   │
│        │                                            │                   │
│        ▼                                            │                   │
│   ┌─────────────────────────────────────────────┐   │                   │
│   │  Pointwise Conv 1×1                         │   │                   │
│   │  Cross-channel feature mixing               │   │                   │
│   └─────────────────────────────────────────────┘   │                   │
│        │                                            │                   │
│        ▼                                            │                   │
│   BatchNorm → GELU                                  │                   │
│        │                                            │                   │
│        ▼                                            │                   │
│   ┌─────────────────────────────────────────────┐   │                   │
│   │          FREQUENCY GATE (Channel Attention) │   │                   │
│   │  ┌───────────────────────────────────────┐  │   │                   │
│   │  │ Global Average Pool → [B, 512, 1, 1]  │  │   │                   │
│   │  └───────────────────────────────────────┘  │   │                   │
│   │                    │                        │   │                   │
│   │                    ▼                        │   │                   │
│   │  ┌───────────────────────────────────────┐  │   │                   │
│   │  │ Conv 1×1 (512→32) → GELU              │  │   │                   │
│   │  │ Conv 1×1 (32→512) → Sigmoid           │  │   │                   │
│   │  │ (Squeeze-Excitation pattern)          │  │   │                   │
│   │  └───────────────────────────────────────┘  │   │                   │
│   │                    │                        │   │                   │
│   │                    ▼                        │   │                   │
│   │           Attention Weights [B, 512, 1, 1]  │   │                   │
│   └─────────────────────────────────────────────┘   │                   │
│        │                                            │                   │
│        ▼                                            │                   │
│   features × attention (element-wise)               │                   │
│        │                                            │                   │
│        ▼                                            │                   │
│   ┌─────────────────────────────────────────────┐   │                   │
│   │  Second Conv Block (same structure)         │   │                   │
│   │  Depthwise 7×7 → Pointwise → BN → GELU      │   │                   │
│   └─────────────────────────────────────────────┘   │                   │
│        │                                            │                   │
│        ▼                                            │                   │
│        + ◄──────────────────────────────────────────┘                   │
│        │                                                                │
│        ▼                                                                │
│   Output [B, 512, H, W]                                                 │
│                                                                         │
└─────────────────────────────────────────────────────────────────────────┘
\end{verbatim}

\textbf{Design Rationale:} - \textbf{Large 7×7 kernels}: Implicitly
capture frequency-domain patterns without explicit FFT -
\textbf{Depthwise separable}: Reduces parameters from O(C²k²) to O(C·k²
+ C²) - \textbf{Frequency gate}: Learns to emphasize/suppress specific
frequency bands - \textbf{Residual connection}: Enables learning
identity + perturbation

\paragraph{4.3.3 Perturbation
Constraints}\label{perturbation-constraints}

The generator output is carefully constrained:

\begin{enumerate}
\def\labelenumi{\arabic{enumi}.}
\tightlist
\item
  \textbf{Tanh activation}: Bounds raw perturbation to {[}-1, 1{]}
\item
  \textbf{Epsilon scaling}:
  \texttt{scaled\_perturbation\ =\ 0.03\ ×\ perturbation}
\item
  \textbf{Clamping}: Ensures adversarial image stays in valid range
\end{enumerate}

\textbf{Maximum perturbation magnitude:} 0.03 (approximately 7.65/255 in
8-bit)

\begin{center}\rule{0.5\linewidth}{0.5pt}\end{center}

\subsection{5. Training Framework}\label{training-framework}

\subsubsection{5.1 Training Philosophy}\label{training-philosophy}

Our training differs fundamentally from traditional GAN training:

\begin{verbatim}
┌─────────────────────────────────────────────────────────────────┐
│                    TRAINING PHILOSOPHY COMPARISON                │
├────────────────────────┬────────────────────────────────────────┤
│    Traditional GAN     │         Our Approach                   │
├────────────────────────┼────────────────────────────────────────┤
│                        │                                        │
│  Generator creates     │  Generator creates adversarial         │
│  realistic fake images │  perturbations to test discriminator   │
│                        │                                        │
│  Discriminator tries   │  Discriminator learns to be robust     │
│  to distinguish fakes  │  despite adversarial attacks           │
│                        │                                        │
│  Goal: D(G(z)) ≈ 0.5   │  Goal: D(x + δ) still correct          │
│  (equilibrium)         │  (adversarial robustness)              │
│                        │                                        │
│  Final product:        │  Final product:                        │
│  Trained generator     │  Robust discriminator                  │
│                        │                                        │
└────────────────────────┴────────────────────────────────────────┘
\end{verbatim}

\subsubsection{5.2 Data Pipeline}\label{data-pipeline}

\paragraph{5.2.1 Dataset: Celeb-DF v2}\label{dataset-celeb-df-v2}

\begin{longtable}[]{@{}llll@{}}
\toprule\noalign{}
Split & Real Images & Fake Images & Total \\
\midrule\noalign{}
\endhead
\bottomrule\noalign{}
\endlastfoot
Train & \textasciitilde65,000 & \textasciitilde472,000 &
\textasciitilde537,000 \\
Test & \textasciitilde36,000 & \textasciitilde261,000 &
\textasciitilde297,000 \\
\end{longtable}

\textbf{Class Imbalance:} \textasciitilde88\% fake, \textasciitilde12\%
real → Addressed with \texttt{pos\_weight=3.0} in BCE loss

\paragraph{5.2.2 Data Transformations}\label{data-transformations}

\begin{verbatim}
┌─────────────────────────────────────────────────────────────────┐
│                     DATA TRANSFORMATION PIPELINE                 │
├─────────────────────────────────────────────────────────────────┤
│                                                                 │
│   ┌─────────────────────────────────────────────────────────┐   │
│   │                  TRAINING TRANSFORMS                    │   │
│   ├─────────────────────────────────────────────────────────┤   │
│   │  1. RandomHorizontalFlip(p=0.5)                         │   │
│   │  2. ColorJitter(                                        │   │
│   │       brightness=0.2,                                   │   │
│   │       contrast=0.2,                                     │   │
│   │       saturation=0.2,                                   │   │
│   │       hue=0.1                                           │   │
│   │     )                                                   │   │
│   │  3. RandomRotation(degrees=10)                          │   │
│   │  4. ToTensor()                                          │   │
│   │  5. Normalize(                                          │   │
│   │       mean=[0.485, 0.456, 0.406],  # ImageNet stats     │   │
│   │       std=[0.229, 0.224, 0.225]                         │   │
│   │     )                                                   │   │
│   └─────────────────────────────────────────────────────────┘   │
│                                                                 │
│   ┌─────────────────────────────────────────────────────────┐   │
│   │                 VALIDATION TRANSFORMS                   │   │
│   ├─────────────────────────────────────────────────────────┤   │
│   │  1. ToTensor()                                          │   │
│   │  2. Normalize(                                          │   │
│   │       mean=[0.485, 0.456, 0.406],                       │   │
│   │       std=[0.229, 0.224, 0.225]                         │   │
│   │     )                                                   │   │
│   └─────────────────────────────────────────────────────────┘   │
│                                                                 │
└─────────────────────────────────────────────────────────────────┘
\end{verbatim}

\paragraph{5.2.3 DataLoader
Configuration}\label{dataloader-configuration}

\begin{Shaded}
\begin{Highlighting}[]
\NormalTok{DataLoader(}
\NormalTok{    batch\_size}\OperatorTok{=}\DecValTok{32}\NormalTok{,}
\NormalTok{    shuffle}\OperatorTok{=}\VariableTok{True}\NormalTok{,           }\CommentTok{\# Random sampling}
\NormalTok{    num\_workers}\OperatorTok{=}\DecValTok{4}\NormalTok{,          }\CommentTok{\# Parallel data loading}
\NormalTok{    pin\_memory}\OperatorTok{=}\VariableTok{True}\NormalTok{,        }\CommentTok{\# Fast CPU→GPU transfer}
\NormalTok{    persistent\_workers}\OperatorTok{=}\VariableTok{True}\NormalTok{, }\CommentTok{\# Keep workers alive}
\NormalTok{    drop\_last}\OperatorTok{=}\VariableTok{True}          \CommentTok{\# Stable batch sizes}
\NormalTok{)}
\end{Highlighting}
\end{Shaded}

\subsubsection{5.3 Training Algorithm}\label{training-algorithm}

\paragraph{5.3.1 Single Training Step}\label{single-training-step}

\begin{verbatim}
┌─────────────────────────────────────────────────────────────────────────┐
│                     TRAINING STEP ALGORITHM                              │
├─────────────────────────────────────────────────────────────────────────┤
│                                                                         │
│   INPUT: Batch of images X with labels Y (0=real, 1=fake)               │
│                                                                         │
│   ═══════════════════════════════════════════════════════════════════   │
│   ║                STEP 1: SEPARATE CLASSES                         ║   │
│   ═══════════════════════════════════════════════════════════════════   │
│                                                                         │
│   X_real = X[Y == 0]    # Real images                                   │
│   X_fake = X[Y == 1]    # Fake images                                   │
│                                                                         │
│   ═══════════════════════════════════════════════════════════════════   │
│   ║             STEP 2: TRAIN DISCRIMINATOR                         ║   │
│   ═══════════════════════════════════════════════════════════════════   │
│                                                                         │
│   ┌─────────────────────────────────────────────────────────────────┐   │
│   │ 2a. Real Image Classification                                   │   │
│   │     D_real = D(X_real)                                          │   │
│   │     L_real = BCE(D_real, 1)   # Should output 1 (real)          │   │
│   └─────────────────────────────────────────────────────────────────┘   │
│                                                                         │
│   ┌─────────────────────────────────────────────────────────────────┐   │
│   │ 2b. Fake Image Classification                                   │   │
│   │     D_fake = D(X_fake)                                          │   │
│   │     L_fake = BCE(D_fake, 0)   # Should output 0 (fake)          │   │
│   └─────────────────────────────────────────────────────────────────┘   │
│                                                                         │
│   ┌─────────────────────────────────────────────────────────────────┐   │
│   │ 2c. Adversarial Robustness (Key Innovation!)                    │   │
│   │     X_adv, δ = G(X_real)      # Generate adversarial images     │   │
│   │     D_adv = D(X_adv.detach()) # Discriminator on adversarial    │   │
│   │     L_adv = BCE(D_adv, 1)     # Should STILL classify as real   │   │
│   └─────────────────────────────────────────────────────────────────┘   │
│                                                                         │
│   ┌─────────────────────────────────────────────────────────────────┐   │
│   │ 2d. Combined Discriminator Loss                                 │   │
│   │     L_D = L_real + L_fake + 0.5 × L_adv                         │   │
│   └─────────────────────────────────────────────────────────────────┘   │
│                                                                         │
│   ┌─────────────────────────────────────────────────────────────────┐   │
│   │ 2e. Update Discriminator                                        │   │
│   │     zero_grad() → backward(L_D) → clip_grad(1.0) → step()       │   │
│   └─────────────────────────────────────────────────────────────────┘   │
│                                                                         │
│   ═══════════════════════════════════════════════════════════════════   │
│   ║               STEP 3: TRAIN GENERATOR                           ║   │
│   ═══════════════════════════════════════════════════════════════════   │
│                                                                         │
│   ┌─────────────────────────────────────────────────────────────────┐   │
│   │ 3a. Generate Fresh Adversarial Images                           │   │
│   │     X_adv, δ = G(X_real)      # New forward pass                │   │
│   └─────────────────────────────────────────────────────────────────┘   │
│                                                                         │
│   ┌─────────────────────────────────────────────────────────────────┐   │
│   │ 3b. Fool Discriminator (Adversarial Loss)                       │   │
│   │     D_adv = D(X_adv)          # NO detach - need gradients      │   │
│   │     L_G_adv = BCE(D_adv, 0)   # Want D to misclassify as fake   │   │
│   └─────────────────────────────────────────────────────────────────┘   │
│                                                                         │
│   ┌─────────────────────────────────────────────────────────────────┐   │
│   │ 3c. Perturbation Regularization                                 │   │
│   │     L_G_perturb = mean(|δ|)   # L1 norm encourages sparsity     │   │
│   └─────────────────────────────────────────────────────────────────┘   │
│                                                                         │
│   ┌─────────────────────────────────────────────────────────────────┐   │
│   │ 3d. Combined Generator Loss                                     │   │
│   │     L_G = L_G_adv + 0.1 × L_G_perturb                           │   │
│   └─────────────────────────────────────────────────────────────────┘   │
│                                                                         │
│   ┌─────────────────────────────────────────────────────────────────┐   │
│   │ 3e. Update Generator                                            │   │
│   │     zero_grad() → backward(L_G) → clip_grad(1.0) → step()       │   │
│   └─────────────────────────────────────────────────────────────────┘   │
│                                                                         │
│   OUTPUT: Metrics (d_loss, g_loss, d_acc_real, d_acc_fake, g_acc_adv)   │
│                                                                         │
└─────────────────────────────────────────────────────────────────────────┘
\end{verbatim}

\subsubsection{5.4 Loss Functions}\label{loss-functions}

\paragraph{5.4.1 Binary Cross-Entropy with
Logits}\label{binary-cross-entropy-with-logits}

\[\mathcal{L}_{BCE} = -\frac{1}{N} \sum_{i=1}^{N} \left[ y_i \cdot \log(\sigma(x_i)) + (1-y_i) \cdot \log(1-\sigma(x_i)) \right]\]

Where \(\sigma(x) = \frac{1}{1+e^{-x}}\) is the sigmoid function.

\textbf{With Class Weighting:}

\[\mathcal{L}_{BCE}^{weighted} = -\frac{1}{N} \sum_{i=1}^{N} \left[ w_{pos} \cdot y_i \cdot \log(\sigma(x_i)) + (1-y_i) \cdot \log(1-\sigma(x_i)) \right]\]

We use \texttt{pos\_weight\ =\ 3.0} to address the 88\%/12\% class
imbalance.

\paragraph{5.4.2 Discriminator Loss}\label{discriminator-loss}

\[\mathcal{L}_D = \mathcal{L}_{real} + \mathcal{L}_{fake} + \lambda_{adv} \cdot \mathcal{L}_{adv}\]

Where: - \(\mathcal{L}_{real} = BCE(D(x_{real}), 1)\) -
\(\mathcal{L}_{fake} = BCE(D(x_{fake}), 0)\) -
\(\mathcal{L}_{adv} = BCE(D(x_{real} + \epsilon \cdot G(x_{real})), 1)\)
- \(\lambda_{adv} = 0.5\) (adversarial weight)

\paragraph{5.4.3 Generator Loss}\label{generator-loss}

\[\mathcal{L}_G = \mathcal{L}_{G,adv} + \lambda_{perturb} \cdot \mathcal{L}_{perturb}\]

Where: -
\(\mathcal{L}_{G,adv} = BCE(D(x_{real} + \epsilon \cdot G(x_{real})), 0)\)
(fool discriminator) - \(\mathcal{L}_{perturb} = \mathbb{E}[|\delta|]\)
(L1 regularization) - \(\lambda_{perturb} = 0.1\) (perturbation weight)

\subsubsection{5.5 Optimization Strategy}\label{optimization-strategy}

\paragraph{5.5.1 Optimizers}\label{optimizers}

\begin{Shaded}
\begin{Highlighting}[]
\CommentTok{\# Discriminator Optimizer}
\NormalTok{AdamW(}
\NormalTok{    params}\OperatorTok{=}\NormalTok{discriminator.parameters(),}
\NormalTok{    lr}\OperatorTok{=}\FloatTok{1e{-}4}\NormalTok{,}
\NormalTok{    betas}\OperatorTok{=}\NormalTok{(}\FloatTok{0.5}\NormalTok{, }\FloatTok{0.999}\NormalTok{),    }\CommentTok{\# Lower β1 for GAN stability}
\NormalTok{    weight\_decay}\OperatorTok{=}\FloatTok{0.01}
\NormalTok{)}

\CommentTok{\# Generator Optimizer  }
\NormalTok{AdamW(}
\NormalTok{    params}\OperatorTok{=}\NormalTok{generator.parameters(),}
\NormalTok{    lr}\OperatorTok{=}\FloatTok{5e{-}5}\NormalTok{,               }\CommentTok{\# Half of discriminator LR}
\NormalTok{    betas}\OperatorTok{=}\NormalTok{(}\FloatTok{0.5}\NormalTok{, }\FloatTok{0.999}\NormalTok{),}
\NormalTok{    weight\_decay}\OperatorTok{=}\FloatTok{0.01}
\NormalTok{)}
\end{Highlighting}
\end{Shaded}

\textbf{Why AdamW?} - Decoupled weight decay (better than Adam) -
Adaptive learning rates per parameter - Standard for transformers and
modern CNNs

\textbf{Why β1=0.5?} - Standard for GAN training - Reduces momentum
oscillations - Improves adversarial training stability

\paragraph{5.5.2 Learning Rate Schedule}\label{learning-rate-schedule}

\begin{Shaded}
\begin{Highlighting}[]
\NormalTok{CosineAnnealingLR(}
\NormalTok{    optimizer,}
\NormalTok{    T\_max}\OperatorTok{=}\DecValTok{50}\NormalTok{,              }\CommentTok{\# Total epochs}
\NormalTok{    eta\_min}\OperatorTok{=}\FloatTok{1e{-}6}           \CommentTok{\# Minimum LR}
\NormalTok{)}
\end{Highlighting}
\end{Shaded}

\[\eta_t = \eta_{min} + \frac{1}{2}(\eta_{max} - \eta_{min})\left(1 + \cos\left(\frac{t \cdot \pi}{T_{max}}\right)\right)\]

\paragraph{5.5.3 Gradient Management}\label{gradient-management}

\begin{Shaded}
\begin{Highlighting}[]
\CommentTok{\# Gradient Clipping (max\_norm=1.0)}
\NormalTok{torch.nn.utils.clip\_grad\_norm\_(}
\NormalTok{    parameters, }
\NormalTok{    max\_norm}\OperatorTok{=}\FloatTok{1.0}\NormalTok{,}
\NormalTok{    norm\_type}\OperatorTok{=}\DecValTok{2}  \CommentTok{\# L2 norm}
\NormalTok{)}
\end{Highlighting}
\end{Shaded}

\textbf{Purpose:} Prevent gradient explosion in adversarial training

\paragraph{5.5.4 Mixed Precision
Training}\label{mixed-precision-training}

\begin{Shaded}
\begin{Highlighting}[]
\NormalTok{precision }\OperatorTok{=} \StringTok{"16{-}mixed"}  \CommentTok{\# Automatic Mixed Precision}
\end{Highlighting}
\end{Shaded}

\textbf{Benefits:} - \textasciitilde50\% memory reduction -
\textasciitilde2-3x speedup on modern GPUs - Maintains numerical
stability via loss scaling

\begin{center}\rule{0.5\linewidth}{0.5pt}\end{center}

\subsection{6. Implementation Details}\label{implementation-details-1}

\subsubsection{6.1 Configuration
Parameters}\label{configuration-parameters}

\begin{verbatim}
┌─────────────────────────────────────────────────────────────────┐
│                    HYPERPARAMETER SUMMARY                        │
├─────────────────────────┬───────────────────────────────────────┤
│ Parameter               │ Value                                 │
├─────────────────────────┼───────────────────────────────────────┤
│ Input Size              │ 224×224×3 RGB                         │
│ Batch Size              │ 32                                    │
│ Max Epochs              │ 50                                    │
│ Learning Rate (D)       │ 1e-4                                  │
│ Learning Rate (G)       │ 5e-5                                  │
│ Adam Betas              │ (0.5, 0.999)                          │
│ Weight Decay            │ 0.01                                  │
│ Gradient Clip           │ 1.0                                   │
│ Epsilon (perturbation)  │ 0.03                                  │
│ Adversarial Weight      │ 0.5                                   │
│ Perturbation Weight     │ 0.1                                   │
│ Pos Weight (BCE)        │ 3.0                                   │
│ Precision               │ 16-mixed (AMP)                        │
│ Strategy                │ DDP (2 GPUs)                          │
├─────────────────────────┼───────────────────────────────────────┤
│ D Backbone              │ ConvNeXt-Tiny (pretrained)            │
│ G Base Channels         │ 64                                    │
│ DCT Size                │ 224×224                               │
└─────────────────────────┴───────────────────────────────────────┘
\end{verbatim}

\subsubsection{6.2 Callbacks and Logging}\label{callbacks-and-logging}

\begin{longtable}[]{@{}ll@{}}
\toprule\noalign{}
Callback & Purpose \\
\midrule\noalign{}
\endhead
\bottomrule\noalign{}
\endlastfoot
\texttt{ModelCheckpoint} & Save top-3 models by validation accuracy \\
\texttt{LearningRateMonitor} & Log LR to TensorBoard \\
\texttt{EarlyStopping} & Stop if val/accuracy doesn't improve for 10
epochs \\
\texttt{TQDMProgressBar} & Clean progress visualization \\
\end{longtable}

\subsubsection{6.3 Hardware Requirements}\label{hardware-requirements}

\begin{longtable}[]{@{}lll@{}}
\toprule\noalign{}
Resource & Minimum & Recommended \\
\midrule\noalign{}
\endhead
\bottomrule\noalign{}
\endlastfoot
GPU VRAM & 8GB & 16GB+ \\
System RAM & 16GB & 32GB+ \\
Storage & 30GB & 50GB+ \\
GPUs & 1 & 2 (DDP) \\
\end{longtable}

\begin{center}\rule{0.5\linewidth}{0.5pt}\end{center}

\subsection{7. Expected Results}\label{expected-results}

\subsubsection{7.1 Training Progression}\label{training-progression}

\begin{verbatim}
Phase 1: Initial Learning (Epochs 0-5)
├── Discriminator learns quickly
│   ├── d_acc_real: 0.5 → 0.85
│   └── d_acc_fake: 0.5 → 0.85
└── Generator struggles (g_loss high)

Phase 2: Adversarial Competition (Epochs 5-20)
├── Generator improves perturbations
├── d_acc_real may drop temporarily (0.85 → 0.70)
└── Losses stabilize

Phase 3: Equilibrium (Epochs 20-50)
├── Both models reach stable state
├── d_acc oscillates around 0.65-0.75
└── Gradual validation improvement
\end{verbatim}

\subsubsection{7.2 Target Metrics}\label{target-metrics}

\begin{longtable}[]{@{}lll@{}}
\toprule\noalign{}
Metric & Target & Description \\
\midrule\noalign{}
\endhead
\bottomrule\noalign{}
\endlastfoot
Accuracy & \textgreater{} 0.80 & Overall classification correctness \\
Precision & \textgreater{} 0.75 & True fakes / Predicted fakes \\
Recall & \textgreater{} 0.75 & Detected fakes / Actual fakes \\
F1 Score & \textgreater{} 0.75 & Harmonic mean of P and R \\
ROC-AUC & \textgreater{} 0.85 & Threshold-independent performance \\
\end{longtable}

\begin{center}\rule{0.5\linewidth}{0.5pt}\end{center}

\subsection{8. Conclusion}\label{conclusion}

This project presents a novel approach to deepfake detection that
combines:

\begin{enumerate}
\def\labelenumi{\arabic{enumi}.}
\tightlist
\item
  \textbf{Frequency-domain analysis} via DCT features that capture
  manipulation artifacts invisible in spatial domain
\item
  \textbf{Adversarial robustness training} that produces a discriminator
  resilient to perturbations
\item
  \textbf{Transfer learning} from ImageNet adapted for single-channel
  DCT features
\item
  \textbf{Modern architectural choices} including ConvNeXt backbone and
  frequency-aware U-Net generator
\end{enumerate}

The resulting system should generalize better to unseen deepfakes
compared to traditional supervised approaches, as the adversarial
training explicitly optimizes for robustness.

\begin{center}\rule{0.5\linewidth}{0.5pt}\end{center}

\subsection{9. References}\label{references}

\begin{enumerate}
\def\labelenumi{\arabic{enumi}.}
\tightlist
\item
  Liu, Z., et al.~(2022). ``A ConvNet for the 2020s.'' CVPR.
\item
  Rossler, A., et al.~(2019). ``FaceForensics++: Learning to Detect
  Manipulated Facial Images.'' ICCV.
\item
  Li, Y., et al.~(2020). ``Celeb-DF: A Large-scale Challenging Dataset
  for DeepFake Forensics.'' CVPR.
\item
  Durall, R., et al.~(2020). ``Unmasking DeepFakes with simple
  Features.'' arXiv.
\item
  Frank, J., et al.~(2020). ``Leveraging Frequency Analysis for Deep
  Fake Image Recognition.'' ICML.
\end{enumerate}

\begin{center}\rule{0.5\linewidth}{0.5pt}\end{center}

\subsection{Appendix A: Complete Architecture
Specifications}\label{appendix-a-complete-architecture-specifications}

\subsubsection{A.1 Discriminator Parameter
Count}\label{a.1-discriminator-parameter-count}

\begin{longtable}[]{@{}ll@{}}
\toprule\noalign{}
Component & Parameters \\
\midrule\noalign{}
\endhead
\bottomrule\noalign{}
\endlastfoot
DCT Extractor & 0 (buffers only) \\
ConvNeXt Stem & 9,312 \\
Stage 1 (3 blocks) & 442,368 \\
Stage 2 (3 blocks) & 1,327,104 \\
Stage 3 (9 blocks) & 11,943,936 \\
Stage 4 (3 blocks) & 14,155,776 \\
Classification Head & 590 \\
\textbf{Total} & \textbf{\textasciitilde27.8M} \\
\end{longtable}

\subsubsection{A.2 Generator Parameter
Count}\label{a.2-generator-parameter-count}

\begin{longtable}[]{@{}ll@{}}
\toprule\noalign{}
Component & Parameters \\
\midrule\noalign{}
\endhead
\bottomrule\noalign{}
\endlastfoot
Encoder Blocks & 4,497,408 \\
Bottleneck & 3,156,480 \\
Decoder Blocks & 5,317,440 \\
Output Layer & 195 \\
\textbf{Total} & \textbf{\textasciitilde13.0M} \\
\end{longtable}

\subsubsection{A.3 Total System
Parameters}\label{a.3-total-system-parameters}

\begin{longtable}[]{@{}ll@{}}
\toprule\noalign{}
Model & Parameters \\
\midrule\noalign{}
\endhead
\bottomrule\noalign{}
\endlastfoot
Discriminator & 27,878,786 \\
Generator & 12,971,523 \\
\textbf{Total} & \textbf{40,850,309} (\textasciitilde41M) \\
\end{longtable}

\begin{center}\rule{0.5\linewidth}{0.5pt}\end{center}

\emph{End of Technical Report}
